\documentclass[11pt]{article}
\usepackage{mathblatt}
% Präambel des Setzers (werkzeuge/setzer-praeambel.tex, wortgleich); dazu \gteil (Muster T2U)
\makeatletter
\setlength{\mbpfbe}{0mm}% keine Punkte (pruefung.md § 5: nur beim Prüfstein)
\newlength{\gEin}\setlength{\gEin}{0pt}
\let\gSchrift\relax
\renewcommand{\pfkreuzzeile}[1]{\par\noindent\hangindent3.2em\hangafter1\hspace*{1.6em}$\square$\ #1\par}
\newcommand{\gkreuz}[1]{$\square$\ #1\hspace{2em}}
\newcommand{\gkreuzzeile}[2]{\par\noindent\hangindent2.2em\hangafter1\mbox{#1\,$\square$}\ #2\par}% Buchstabe, Option: eine je Zeile (neben Skizze)
% links, was man ansieht, rechts, was man tut; Spaltengrenze überall an derselben Stelle (gemeinsam.md „Satz“)
% v0.17: Fußeintrag „Nr: Wert“ – kurz in einer Box (nie getrennt), lang umbrechbar mit der Nummer am Anfang
% (vorher \mbox: lange Lösungen liefen über den Rand, AS40 138 pt); robust, weil \fusshilfe in die .aux schreibt
\NewDocumentCommand{\gfuss}{m m}{\sbox0{#1:~#2}\ifdim\wd0<0.45\textwidth\mbox{#1:~#2}\else\mbox{#1:}\nobreak\ #2\fi}
\newsavebox{\gbox}\newlength{\gLinks}
% v0.10 (Lehrer 13.10.): keine starre Halbierung – die linke Spalte ist so breit wie die Figur (+ 6 mm), die
% Antwort steht direkt daneben; ist die Figur breiter als 0,62 Zeile, steht die Antwort darunter
\newcommand{\gzwei}[2]{\par\smallskip\noindent\sbox{\gbox}{#1}%
  \setlength{\gLinks}{\dimexpr\wd\gbox+6mm\relax}%
  \ifdim\gLinks>0.62\linewidth
    \usebox{\gbox}\par\smallskip\noindent\begin{minipage}{\linewidth}#2\end{minipage}%
  \else
    \begin{minipage}[t]{\gLinks}\vspace{0pt}\usebox{\gbox}\end{minipage}%
    \begin{minipage}[t]{\dimexpr\linewidth-\gLinks\relax}\vspace{0pt}\raggedright #2\end{minipage}%
  \fi\par}
% 16.10. (Lehrer, PA10 S. 1): Text links, Koordinatensystem rechts daneben; ist das System breiter als 0,58 Zeile,
% steht es wie bisher unter dem Text
% Passt das System knapp nicht mehr auf die Seite (bis 25 % zu hoch), wird es auf den Rest verkleinert statt
% die Seite halb leer zu lassen
\newlength{\gRest}\newlength{\gHoch}
\newcommand{\gtextks}[2]{\sbox{\gbox}{#1}%
  \setlength{\gHoch}{\dimexpr\ht\gbox+\dp\gbox\relax}%
  \ifdim\pagegoal<\maxdimen\setlength{\gRest}{\dimexpr\pagegoal-\pagetotal-4\baselineskip\relax}%
    \ifdim\gHoch>\gRest\ifdim\gRest>0.75\gHoch
      \sbox{\gbox}{\resizebox*{!}{\gRest}{\usebox{\gbox}}}\fi\fi\fi
  \ifdim\wd\gbox>0.58\linewidth #2\par\smallskip\noindent\gpassend{\usebox{\gbox}}\par
  \else\begin{minipage}[t]{\dimexpr\linewidth-\wd\gbox-5mm\relax}\raggedright #2\end{minipage}\hfill
    \begin{minipage}[t]{\wd\gbox}\vspace{0pt}\usebox{\gbox}\end{minipage}\par\fi}
% v0.9 (12.10., pruefung.md § 5 „Marke mit B, Bedingungen rechts“): drittes Argument von pfaufg = Bedingungen
% der Prüfung (v0.10: \pfohnetr = TR durchgestrichen, „FOR“), klein grau im rechten Rand auf der Grundlinie der Nummer, Zeilen mit \\
\newcommand{\gbedrand}[2]{\ifblank{#2}{}{\rlap{\hspace*{\dimexpr#1+2mm\relax}%
  \parbox[t]{15mm}{\raggedright\scriptsize\leavevmode\color{mbgrau}\strut#2}}}}
% Teil eines Bündels (a, b …): Marke links im Rand, Bedingungen rechts im Rand, auf der Zeile des Teils
\newcommand{\gteilrand}[2]{\llap{\mbpfrandmarke{#1}\hspace*{\mbpfnr}}\gbedrand{\linewidth}{#2}}
% Teil eines Bündels: Buchstabe in eigener Spalte, Inhalt eingerückt (\gEin = Einzug, damit die Spaltengrenze von
% \gzwei an derselben Stelle bleibt)
\newenvironment{gbteil}[3]{\par\smallskip\noindent\gteilrand{#2}{#3}\makebox[\mbpfteil][l]{\textbf{#1}}%
  \begin{minipage}[t]{\dimexpr\linewidth-\mbpfteil\relax}\setlength{\gEin}{\mbpfteil}\raggedright\strut\ignorespaces}%
 {\end{minipage}\par\addvspace{6pt}}% v0.18 (PG80 Nr. 1): Abstand, damit die Linie von a) nicht an b) rutscht
% Grafik höchstens so breit wie die Spalte (Bündel nebeneinander, Koordinatensystem ohne Karo)
\newsavebox{\gpbox}
\newcommand{\gpassend}[1]{\sbox{\gpbox}{#1}\ifdim\wd\gpbox>\linewidth\typeout{GPASSEND: verkleinert \the\wd\gpbox\space auf \the\linewidth}\resizebox{\linewidth}{!}{\usebox{\gpbox}}\else\usebox{\gpbox}\fi}
% Skizzenfeld (v0.10, Lehrer 13.10.): Karo ab dem linken Rand über die Zeile, „Skizze“ klein grau oben links im Feld
% (wie \pfskizzenfeld der Vorlage); der rechte Rand bleibt den Bedingungen
\newcommand{\gskizze}[1]{\par\smallskip\noindent\begin{tikzpicture}
  \clip (0,0) rectangle ({\linewidth-1mm},{#1*10mm});
  \draw[black!16,line width=0.3pt,step=5mm] (0,0) grid ({\linewidth},{#1*10mm});
  \node[anchor=north west,inner sep=2pt,font=\scriptsize,text=mbgrau,fill=white] at (0,{#1*10mm}) {Skizze};
  \end{tikzpicture}\par}
\RenewDocumentEnvironment{pfaufg}{m m m}{%
  \begin{lrbox}{\mbaufgabebox}\begin{minipage}[t]{\linewidth}%
  \gSchrift\noindent\hspace*{\gEin}\mbpfrandmarke{#1}\gbedrand{\linewidth-\gEin}{#3}%
  \begin{minipage}[t]{\mbpfnr}\strut\textbf{#2}\end{minipage}%
  \begin{minipage}[t]{\dimexpr\linewidth-\gEin-\mbpfnr\relax}\raggedright\strut\ignorespaces}%
 {\end{minipage}%
  \end{minipage}\end{lrbox}%
  \setlength{\mbaufgabehoehe}{\dimexpr\ht\mbaufgabebox+\dp\mbaufgabebox\relax}%
  \ifdefempty{\gkopfnach}{\par\addvspace{5pt}\Needspace*{\mbaufgabehoehe}}{%
    % Stufenkopf, der mit dieser (ersten) Aufgabe zusammenbleibt: Platz für beide verlangen, dann den Kopf setzen
    % (12.10., Quadratische: ein Kopf stand allein unten, die Aufgabe mit Koordinatensystem auf der nächsten Seite)
    % v0.17 (AS40 „Terme deuten“ allein unten): hinter dem Kopf steht das Sprungziel (\special, ein Knoten, nach
    % dem die Seite brechen durfte) – darum \nobreak nach dem Kopf; Vorhalt 56 pt = Kopf mit Linie und Abständen
    \par\Needspace*{\dimexpr\mbaufgabehoehe+56pt\relax}\gkopfnach\par\nobreak\gdef\gkopfnach{}\addvspace{5pt}}%
  \noindent\usebox{\mbaufgabebox}\par\addvspace{5pt}}
% v0.18 (plan.md § 10, 15.10. l): Abschnittsbuchstabe A, B, C vor dem Namen im Kopf – als Vorsatz, damit der
% Name im TeX gleich bleibt (Bestellkarte, Sprungziele lesen ihn)
\def\gbstakt{}\newcommand{\gbstvor}[1]{\gdef\gbstakt{#1}}
\NewCommandCopy\gpfstufekopfalt\pfstufekopf
\RenewDocumentCommand{\pfstufekopf}{o m m}{\IfValueTF{#1}{\gpfstufekopfalt[#1]}{\gpfstufekopfalt}%
  {\ifx\gbstakt\empty\else\gbstakt\enspace\fi #2}{#3}\gdef\gbstakt{}}
% v0.18 (16.10. e): Begriff grau, ohne Nummer, ohne Rahmen über der ersten Aufgabe der Stufe
\newcommand{\gbegriff}[1]{\par\nobreak\noindent\hspace*{\mbpfnr}\parbox[t]{\dimexpr\linewidth-\mbpfnr\relax}{\raggedright\small\color{mbgrau}#1}\par\nobreak\addvspace{3pt}}
% v0.18 (16.10. g): Lösungsblatt – Kopf mit kräftigem Strich, Abschnitt klein grau als Zwischenzeile
\newcommand{\glkopf}[1]{\noindent{\large\bfseries Lösungen \textperiodcentered{} #1}\par\vspace{3pt}\hrule height 1.2pt\par\medskip}
\newcommand{\glzwischen}[1]{\par\addvspace{8pt}\Needspace*{4\baselineskip}\noindent{\footnotesize\color{mbgrau}#1}\par\nobreak\vspace{1pt}}
% v0.18 (15.10. o, w): Blattende – Kern ohne Rahmen, dann graue Zeile „Mehr: n Aufgaben“
\newcommand{\gkern}[1]{\par\addvspace{14pt}\Needspace*{6\baselineskip}\noindent{\bfseries Kern}\par\nobreak\smallskip
  {\small\raggedright #1\par}}
\newcommand{\gmehr}[1]{\par\addvspace{10pt}\noindent{\footnotesize\color{mbgrau}Mehr: #1 Aufgaben}\par}
% Stufenkopf vormerken: \gkopfvor{<Kopf>} – die nächste pfaufg setzt ihn, sobald beide zusammen Platz haben
\def\gkopfnach{}
\newcommand{\gkopfvor}[1]{\gdef\gkopfnach{#1}}
% Teilaufgabe mit Skizze links in gzwei: Buchstabe in eigener Spalte, darunter Text und Skizze
\newcommand{\gteil}[2]{\makebox[\mbpfteil][l]{\textbf{#1}}\parbox[t]{\dimexpr0.5\textwidth-\gEin-\mbpfnr-5mm-\mbpfteil\relax}{\raggedright #2}}
% Stufenkopf der Vorlage (mathblatt.sty 2026-10-12: Linie, Abstand, 13 pt, \pfjahre) ohne ihren Vorhalt von
% 14 Zeilen: der Setzer fordert vor dem Kopf 0,16 Seitenhöhe (Kopf + erste Aufgabe), mehr nicht
\setlength{\mbpfstufevorhalt}{0pt}
\makeatother
\newcommand{\karo}[1]{\par\smallskip\noindent\begin{tikzpicture}
  \clip (0,0) rectangle ({\linewidth-1mm},{#1*10mm});
  \draw[black!16,line width=0.3pt,step=5mm] (0,0) grid ({\linewidth},{#1*10mm});
  \end{tikzpicture}\par}
\newcommand{\kk}[1]{$\square$\,#1\hspace{0.9em}}
\newcommand{\linie}[1]{\,{\color{black!45}\rule[-1pt]{#1}{0.4pt}}\,}
\newcommand{\pkt}{\textperiodcentered{}}
% v0.14 (14.10.): Stufenkopf mit fünf Punkten (Option = Jahre für die Bestellkarte) oder „zuletzt 19“
\NewDocumentCommand{\gpunkte}{o m}{#2}
\newcommand{\gzuletzt}[1]{zuletzt #1}
% v0.7 (12.10.): Schülerfrage unter dem Stufenkopf; Ankreuzoption als Bild (Buchstabe und Kästchen darunter)
\newcommand{\gfrage}[1]{\par\nobreak\noindent\hspace*{\mbpfnr}{\small\itshape #1}\par\nobreak\addvspace{2pt}}
\newlength{\gbildbreite}\setlength{\gbildbreite}{0.32\linewidth}% Setzer setzt sie je Reihe (vier Bilder: 0,24)
\newcommand{\gbildoption}[3]{\begin{minipage}[t][#3][b]{\gbildbreite}\centering #2\par\smallskip\mbox{#1\,$\square$}\end{minipage}}
% Lineare (12.10.): Bild mit Buchstabe ohne Kästchen (Diagramm zuordnen); Aussage mit „wahr / falsch“ zum Ankreuzen
\newcommand{\gbildname}[3]{\begin{minipage}[t][#3][b]{\gbildbreite}\centering #2\par\smallskip\mbox{\textbf{#1}}\end{minipage}}
\newcommand{\gwahrfalsch}[2]{\par\noindent\hangindent2.2em\hangafter1\makebox[2.2em][l]{(#1)}#2\hfill\mbox{wahr $\square$\quad falsch $\square$}\par}
\usepackage{graphicx}\graphicspath{{/home/claude/agent-d1/aufgabenbank-privat/}}
\begin{document}

\pfheftstil
\pflaufkopf[P10 \pkt{} Prüfung]{Pythagoras}{PG50-6}\pfseitevon
\pfheftkopf{Pythagoras}{}
\pfstufenzeile{A~Pythagoras aufstellen \pkt{} B~Hypotenuse gesucht \pkt{} C~Kathete gesucht \pkt{} D~Pythagoras versteckt}
% Formelkasten (15.10. t, 16.10. a)
{\small% duden-s21-rechtwinkliges-dreieck – Rechtwinkliges Dreieck (Duden S. 21), aus duden-sek1.md
\providecommand\fsf[1]{\leavevmode\setbox0\hbox{\everymath{\displaystyle}#1}\ifdim\wd0>\linewidth\setbox0\hbox{\resizebox{\linewidth}{!}{\box0}}\fi
 \vrule height\dimexpr\ht0+1.5pt\relax depth\dimexpr\dp0+1.5pt\relax width0pt\box0}
\providecommand\fstab[1]{\leavevmode\setbox0\hbox{#1}\ifdim\wd0>\linewidth\resizebox{\linewidth}{!}{\box0}\else\box0\fi}
\providecommand\fssic[1]{{\color{red!65!black}\footnotesize[sic?#1]}}
\providecommand\fsunl[1]{{\color{red!65!black}\footnotesize[unleserlich#1]}}
\providecommand\fsfigur[1]{{\footnotesize\color{mbgrau}\itshape Figur im Original: #1\par}}
\providecommand\fsbild[1]{\leavevmode\hfill#1\par}
\providecommand\fsbildbreite{4.6cm}
\providecommand\fszeilebild[6]{\par\noindent\begin{minipage}[t]{\dimexpr\fsbreite-#6-6pt}\vspace{0pt}\renewcommand\fsbreite{\linewidth}\fszeile{#1}{#2}{#3}{}\end{minipage}\hfill
 \begin{minipage}[t]{#6}\vspace{3pt}\includegraphics[width=#6]{#5}\end{minipage}\par}
\usetikzlibrary{angles}\tikzset{fsfig/.style={line width=0.7pt,baseline=(current bounding box.north),font=\small},
 fswinkel/.style={draw=mbgrau,line width=0.5pt,angle radius=5.5mm,angle eccentricity=0.62,pic text options={font=\footnotesize}},
 fsrecht/.style={draw=mbgrau,line width=0.5pt,angle radius=2.6mm,angle eccentricity=0.5,pic text={\tiny$\bullet$}},
 fslinie/.style={font=\scriptsize,sloped,inner sep=1.5pt},fsflaeche/.style={fill=mbgrau,fill opacity=0.18,draw=none}}
\providecommand\fspkt[1]{\fill #1 circle (1.4pt);}
\providecommand\fshinweis[1]{{\footnotesize\color{mbgrau}#1\par}}
\providecommand\fspunkt{\ensuremath{\scriptstyle\blacksquare}\,}
\providecommand\fsfuenf{\ensuremath{\rlap{$||||$}\diagup}}
\providecommand\fsquelle[1]{\hfill{\footnotesize\color{mbgrau}#1}}
\providecommand\fsbreite{\linewidth}
\providecommand\fszeile[4]{\par\noindent{\renewcommand\arraystretch{1.25}\setlength\lineskip{4pt}\setlength\lineskiplimit{3pt}%
 \ifblank{#2}%
 {\begin{tabular}{>{\columncolor{mbkasten}\raggedright\hspace{0pt}}p{\dimexpr0.25\fsbreite-2\tabcolsep}>{\raggedright\arraybackslash}p{\dimexpr0.75\fsbreite-2\tabcolsep}}%
  #1 & #3\ifblank{#4}{}{\\ & #4}\end{tabular}}%
 {\begin{tabular}{>{\columncolor{mbkasten}\raggedright\hspace{0pt}}p{\dimexpr0.25\fsbreite-2\tabcolsep}>{\raggedright}p{\dimexpr0.33\fsbreite-2\tabcolsep}>{\raggedright\arraybackslash}p{\dimexpr0.42\fsbreite-2\tabcolsep}}%
  #1 & #2 & #3\ifblank{#4}{}{\\ & \multicolumn{2}{>{\raggedright\arraybackslash}p{\dimexpr0.75\fsbreite-2\tabcolsep}}{#4}}\end{tabular}}}\par}
\providecommand\fskastentitel[1]{\par\noindent\textbf{#1}\par}
\uebersichtskasten{\renewcommand\fsbreite{\linewidth}
\fszeilebild{\textbf{Rechtwinkliges Dreieck}}{\fspunkt{} eine Hypotenuse und zwei Katheten\newline \fspunkt{} Hypotenuse (Seite gegenüber dem rechten Winkel) ist längste Seite}{\fsf{$c^2 = a^2 + b^2$ Satz des PYTHAGORAS}\newline \fsf{$h^2 = q\cdot p$ Höhensatz}\newline \fsf{$a^2 = c\cdot p$ Kathetensatz}\newline \fsf{$b^2 = c\cdot q$ Kathetensatz}\newline \fsf{$u = a+b+c$}\newline \fsf{$A = \frac{a\cdot b}{2}$}\newline a, b Katheten\newline c Hypotenuse}{}{formelsammlung/figuren/duden-bild/duden-s21-rechtwinkliges-dreieck.png}{3.65cm}
}
}\vspace{-4pt}\par\noindent{\footnotesize\color{mbgrau}Formeln und Werte S.~21}\par\medskip

% ===================== Stufe: Gleichung aufstellen =====================
\Needspace*{0.16\textheight}\gbstvor{A}\pfstufekopf{Pythagoras aufstellen}{\gpunkte[26 24 22]{\textcolor{black!80}{$\bullet$}\,\textcolor{black!25}{$\circ$}\,\textcolor{black!80}{$\bullet$}\,\textcolor{black!25}{$\circ$}\,\textcolor{black!80}{$\bullet$}}}\special{pdf:dest (stufe-pythagoras-aufstellen) [@thispage /XYZ @xpos @ypos null]}
% Schülerfrage der Stufe (Gliederung; Einstieg nur von dort, nie aus der Bank – 12.10.)
\begin{pfaufg}{}{1.}{}
Markiere die Hypotenuse. Schreib dann die Gleichung des Satzes des Pythagoras auf.
\gzwei{\begin{tikzpicture}[line width=0.7pt,scale=0.8]\coordinate (P1) at (0,0);\coordinate (P2) at (2.5,0);\coordinate (P3) at (0.9,1.2);\draw (P1) -- (P2) -- (P3) -- cycle;\pfrw[0.35]{(P3)}{-126.9}\node[font=\small,anchor=north] at (1.25,-0.08) {$r$};\node[font=\small,anchor=south east] at (0.40,0.62) {$p$};\node[font=\small,anchor=south west] at (1.75,0.62) {$q$};\end{tikzpicture}}{\vspace{-\baselineskip}\par\noindent \hfill\linie{30mm}}
\par\fusshilfe{\gfuss{1}{r² = p² + q²}}
\end{pfaufg}
% 2022-OS-B1g-m1 (Original 2022-OS-B1g, Schwere I, rang 1; herausgelöste Fassung: 2022-OS-B1g-h1)
\begin{pfaufg}{’22~\pkt{}~1g}{2.}{Basis}
Kreuze die Aussage an, die für jedes rechtwinklige Dreieck stimmt:
\par\smallskip\noindent\pfkreuz{}{Die Seite gegenüber dem rechten Winkel ist eine Kathete.}\pfkreuz{}{$c=a+b$}\pfkreuz{}{Die beiden Kathetenquadrate haben zusammen denselben Flächeninhalt wie das Hypotenusenquadrat.}
\fusshilfe{\gfuss{2}{Die beiden Kathetenquadrate haben zusammen denselben Flächeninhalt wie das Hypotenusenquadrat}}
\end{pfaufg}
% 2026-FOR-B1j-m1 (Original 2026-FOR-B1j, Schwere I, rang 2; herausgelöste Fassung: 2026-FOR-B1j-h1)
\begin{pfaufg}{’26~\pkt{}~1j}{3.}{Basis}
Im abgebildeten Dreieck sind u und v die Katheten, w ist die Hypotenuse. Kreuze an, mit welcher Gleichung man w berechnen kann:
\gzwei{\begin{tikzpicture}[line width=0.7pt,font=\small,scale=0.85]\coordinate (P1) at (0,0);\coordinate (P2) at (2.8,0);\coordinate (P3) at (2.8,1.7);\draw (P1) -- (P2) -- (P3) -- cycle;\pfrw[0.35]{(P2)}{90}\node[anchor=north] at (1.4,-0.08) {$u$};\node[anchor=west] at (2.88,0.85) {$v$};\node[anchor=south east] at (1.36,0.9) {$w$};\end{tikzpicture}}{\gkreuzzeile{}{$w=\sqrt{v^2-u^2}$}\gkreuzzeile{}{$w=\sqrt{u^2-v^2}$}\gkreuzzeile{}{$w=\sqrt{u^2+v^2}$}}
\fusshilfe{\gfuss{3}{$w=\sqrt{u^2+v^2}$}}
\end{pfaufg}
% 2024-OS-B1f-m1 (Original 2024-OS-B1f, Schwere I, rang 3; herausgelöste Fassung: 2024-OS-B1f-h1)
\begin{pfaufg}{’24~\pkt{}~1f}{4.}{Basis}
Im abgebildeten Dreieck sind y und z die Katheten, x ist die Hypotenuse. Kreuze die Gleichung an, mit der man die Länge von z berechnen kann:
\gzwei{\begin{tikzpicture}[line width=0.7pt,font=\small,scale=0.85]\coordinate (P1) at (0,0);\coordinate (P2) at (2.8,0);\coordinate (P3) at (0,1.7);\draw (P1) -- (P2) -- (P3) -- cycle;\pfrw[0.35]{(P1)}{0}\node[anchor=north] at (1.4,-0.08) {$z$};\node[anchor=east] at (-0.08,0.85) {$y$};\node[anchor=south west] at (1.44,0.9) {$x$};\end{tikzpicture}}{\gkreuzzeile{}{$z=x^2-y^2$}\gkreuzzeile{}{$z=\sqrt{x^2+y^2}$}\gkreuzzeile{}{$z=x+y$}\gkreuzzeile{}{$z=\sqrt{x^2-y^2}$}}
\fusshilfe{\gfuss{4}{$z=\sqrt{x^2-y^2}$}}
\end{pfaufg}

% ===================== Stufe: Hypotenuse gesucht =====================
\Needspace*{0.16\textheight}\gbstvor{B}\pfstufekopf{Hypotenuse gesucht}{\gpunkte[25]{\textcolor{black!25}{$\circ$}\,\textcolor{black!80}{$\bullet$}\,\textcolor{black!25}{$\circ$}\,\textcolor{black!25}{$\circ$}\,\textcolor{black!25}{$\circ$}}}\special{pdf:dest (stufe-hypotenuse-gesucht) [@thispage /XYZ @xpos @ypos null]}
% Schülerfrage der Stufe (Gliederung; Einstieg nur von dort, nie aus der Bank – 12.10.)
\begin{pfaufg}{}{5.}{}
Bei welcher Seite ist die Hypotenuse gesucht?
\gzwei{\begin{tikzpicture}[line width=0.7pt,scale=0.7,font=\footnotesize]\draw (0,0)--(2.4,0)--(2.4,1.6)--cycle;\pfrw[0.25]{(2.4,0)}{90}\node[anchor=north] at (1.2,-0.05) {$6$ cm};\node[anchor=west] at (2.45,0.8) {$4$ cm};\node[anchor=south east] at (1.2,0.85) {$x$};\begin{scope}[xshift=4.8cm]\draw (0,0)--(2.4,0)--(0,1.6)--cycle;\pfrw[0.25]{(0,0)}{0}\node[anchor=north] at (1.2,-0.05) {$y$};\node[anchor=east] at (-0.05,0.8) {$4$ cm};\node[anchor=south west] at (1.25,0.85) {$7$ cm};\end{scope}\begin{scope}[xshift=9.6cm]\draw (0,0.3)--(0.99,1.34)--(2.4,0)--cycle;\pfrw[0.25]{(0.99,1.34)}{-133.6}\node[anchor=south east] at (0.45,0.85) {$5$ cm};\node[anchor=south west] at (1.75,0.75) {$6$ cm};\node[anchor=north] at (1.2,0.1) {$z$};\end{scope}\end{tikzpicture}}{\vspace{-\baselineskip}\par\noindent\makebox[\mbpfteil][l]{\textbf{a)}}x\hfill\kk{ja}\par\noindent\makebox[\mbpfteil][l]{\textbf{b)}}y\hfill\kk{ja}\par\noindent\makebox[\mbpfteil][l]{\textbf{c)}}z\hfill\kk{ja}}
\par\fusshilfe{\gfuss{5}{a) und c)}}
\end{pfaufg}
% 2025-OS-K4a-m1 (Original 2025-OS-K4a, Schwere I, rang 4; herausgelöste Fassung: 2025-OS-K4a-h2)
% Zweitthema ausgeblendet (Spalten <feld>_einheit, pruefheft-p.md v0.9)
\begin{pfaufg}{nach~P10\\’25~\pkt{}~4a}{6.}{}
Familie Yücel baut an die Stufe vor der Haustür eine feste Rampe für Rollstühle. Die Stufe ist 17 cm hoch; waagerecht ist die Rampe 200 cm lang. Berechne die Länge x der Rampe.
\gzwei{\begin{tikzpicture}[line width=0.7pt,font=\small,scale=0.75]\coordinate (P1) at (0,0);\coordinate (P2) at (6.0,0);\coordinate (P3) at (6.0,0.9);\draw (P1) -- (P2) -- (P3) -- cycle;\pfrw[0.3]{(P2)}{90}\node[anchor=north] at (3.0,-0.08) {$200$ cm};\node[anchor=west] at (6.08,0.45) {$17$ cm};\node[anchor=south east] at (3.0,0.5) {$x$};\end{tikzpicture}}{\karo{2}}
\fusshilfe{\gfuss{6}{200,7 cm}}
\end{pfaufg}

% ===================== Stufe: Kathete gesucht =====================
\Needspace*{0.16\textheight}\gbstvor{C}\pfstufekopf{Kathete gesucht}{\gpunkte[26 24 22]{\textcolor{black!80}{$\bullet$}\,\textcolor{black!25}{$\circ$}\,\textcolor{black!80}{$\bullet$}\,\textcolor{black!25}{$\circ$}\,\textcolor{black!80}{$\bullet$}}}\special{pdf:dest (stufe-kathete-gesucht) [@thispage /XYZ @xpos @ypos null]}
% Schülerfrage der Stufe (Gliederung; Einstieg nur von dort, nie aus der Bank – 12.10.)
\begin{pfaufg}{}{7.}{}
Kreuze an, mit welcher Gleichung du die Länge der Seite k berechnest.
\gzwei{\begin{tikzpicture}[line width=0.7pt,scale=0.85]\coordinate (P1) at (0,0);\coordinate (P2) at (2.8,0);\coordinate (P3) at (0,1.7);\draw (P1) -- (P2) -- (P3) -- cycle;\pfrw[0.35]{(P1)}{0}\node[font=\small,anchor=north] at (1.4,-0.08) {$g$};\node[font=\small,anchor=east] at (-0.08,0.85) {$k$};\node[font=\small,anchor=south west] at (1.44,0.9) {$m$};\end{tikzpicture}}{\vspace{-\baselineskip}\par\noindent\makebox[\mbpfteil][l]{\textbf{a)}}$k = \sqrt{m^2 + g^2}$\hfill\kk{ja}\par\noindent\makebox[\mbpfteil][l]{\textbf{b)}}$k = \sqrt{m^2 - g^2}$\hfill\kk{ja}\par\noindent\makebox[\mbpfteil][l]{\textbf{c)}}$k = \sqrt{g^2 - m^2}$\hfill\kk{ja}}
\par\fusshilfe{\gfuss{7}{b)}}
\end{pfaufg}
% 2026-FOR-K4a-m1 (Original 2026-FOR-K4a, Schwere I, rang 5; herausgelöste Fassung: 2026-FOR-K4a-h1)
\begin{pfaufg}{nach~P10\\’26~\pkt{}~4a}{8.}{}
Im Viereck ABCD sind gegenüberliegende Seiten parallel. Die Seite BC ist 38 cm lang. Die Seite AB wird über B hinaus bis F verlängert; BF ist 16 cm lang. Die Höhe h = CF steht senkrecht auf AF. Der Winkel bei B zwischen BF und BC heißt $\varepsilon$. Berechne die Höhe h.
\gzwei{\begin{tikzpicture}[line width=0.7pt,font=\small,scale=0.06]\coordinate (A) at (0,0);\coordinate (B) at (50,0);\coordinate (F) at (66,0);\coordinate (C) at (66,34.5);\coordinate (D) at (16,34.5);\draw (A) -- (B) -- (C) -- (D) -- cycle;\draw (B) -- (F);\draw[dashed] (F) -- (C);\pfrw[5]{(F)}{90}\node[anchor=north east] at (A) {$A$};\node[anchor=north east] at (B) {$B$};\node[anchor=north west] at (F) {$F$};\node[anchor=south west] at (C) {$C$};\node[anchor=south east] at (D) {$D$};\node[anchor=south east] at (57.5,17.2) {$38$ cm};\node[anchor=north] at (60,-1) {$16$ cm};\node[anchor=west] at (67,17.2) {$h$};\end{tikzpicture}}{\karo{2}}
\fusshilfe{\gfuss{8}{34,47 cm}}
\end{pfaufg}
% 2024-OS-K6a-m1 (Original 2024-OS-K6a, Schwere I, rang 6; herausgelöste Fassung: 2024-OS-K6a-h1)
\begin{pfaufg}{nach~P10\\’24~\pkt{}~6a}{9.}{}
Eine Seilbahn führt vom Ort B im Tal zum Ort A auf einem Berg. F liegt senkrecht unter A auf der Höhe von B; bei F ist ein rechter Winkel. FB ist 262 m lang, die Seilbahnstrecke AB 396 m. Eine zweite Seilbahn verbindet B mit dem Ort C auf einem anderen Berg. Bei B heißt der Winkel zwischen BF und BA $\beta_1$, der Winkel zwischen BA und BC $\beta_2$; bei A heißt der Winkel zwischen AB und AC $\alpha$. Die Strecke FA ist der Höhenunterschied zwischen B und A. Berechne diesen Höhenunterschied.
\gzwei{\begin{tikzpicture}[line width=0.7pt,font=\small,scale=0.8]\coordinate (A) at (0,2.97);\coordinate (F) at (0,0);\coordinate (B) at (2.62,0);\draw (A) -- (F) -- (B) -- cycle;\pfrw[0.3]{(F)}{0}\node[anchor=south] at (A) {$A$};\node[anchor=north east] at (F) {$F$};\node[anchor=north west] at (B) {$B$};\node[anchor=north] at (1.31,-0.08) {$262$ m};\node[anchor=south west] at (1.36,1.52) {$396$ m};\end{tikzpicture}}{\karo{2}}
\fusshilfe{\gfuss{9}{296,9 m}}
\end{pfaufg}
% 2022-OS-K5a-m1 (Original 2022-OS-K5a, Schwere I, rang 7; herausgelöste Fassung: 2022-OS-K5a-h1)
\begin{pfaufg}{nach~P10\\’22~\pkt{}~5a}{10.}{}
Gegeben ist ein Dreieck ABC, das nicht rechtwinklig ist. Die Höhe $h_c$ von C auf die Seite AB ist etwa 7,3 m lang; ihr Fußpunkt heißt D. Die Seite b = AC ist 13,2 m lang, der Winkel bei B ist 49° groß. Der Winkel bei C zwischen AC und der Höhe heißt $\gamma_1$. Weise nach, dass die Strecke AD etwa 11,0 m lang ist.
\gzwei{\begin{tikzpicture}[line width=0.7pt,font=\small,scale=0.3]\coordinate (A) at (0,0);\coordinate (B) at (17.4,0);\coordinate (C) at (11,7.3);\coordinate (D) at (11,0);\draw (A) -- (B) -- (C) -- cycle;\draw[dashed] (C) -- (D);\pfrw[0.9]{(D)}{90}\node[anchor=north east] at (A) {$A$};\node[anchor=north west] at (B) {$B$};\node[anchor=south] at (C) {$C$};\node[anchor=north] at (D) {$D$};\node[anchor=south east] at (5.5,3.7) {$13{,}2$ m};\node[anchor=west] at (11.8,3.2) {$h_c = 7{,}3$ m};\end{tikzpicture}}{\karo{2}}
\fusshilfe{\gfuss{10}{11,0 m}}
\end{pfaufg}

% ===================== Stufe: Dreieck erst in Figur oder Körper finden =====================
\Needspace*{0.16\textheight}\gbstvor{D}\pfstufekopf{Pythagoras versteckt}{\gpunkte[26 25 22]{\textcolor{black!80}{$\bullet$}\,\textcolor{black!80}{$\bullet$}\,\textcolor{black!25}{$\circ$}\,\textcolor{black!25}{$\circ$}\,\textcolor{black!80}{$\bullet$}}}\special{pdf:dest (stufe-pythagoras-versteckt) [@thispage /XYZ @xpos @ypos null]}
% Schülerfrage der Stufe (Gliederung; Einstieg nur von dort, nie aus der Bank – 12.10.)
\begin{pfaufg}{}{11.}{}
Welche Strecken bilden das rechtwinklige Dreieck? Nenne die beiden Katheten und die Hypotenuse.
\par\smallskip
\par\noindent\makebox[\mbpfteil][l]{\textbf{a)}}Kegel mit Radius r, Höhe h, Mantellinie s\hfill\linie{30mm}
\par\noindent\makebox[\mbpfteil][l]{\textbf{b)}}Zylinder mit Radius r, Höhe h; Stab vom Rand des Bodens schräg zum gegenüberliegenden oberen Rand\hfill\linie{30mm}
\par\fusshilfe{\gfuss{11}{a) Katheten r und h, Hypotenuse s; b) Katheten 2 · r und h, Hypotenuse der Stab}}
\end{pfaufg}
% 2026-FOR-K2c-m1 (Original 2026-FOR-K2c, Schwere II, rang 8; herausgelöste Fassung: keine)
\begin{pfaufg}{nach~P10\\’26~\pkt{}~2c}{12.}{FOR}
Ein Turm besteht aus einem Zylinder (Höhe 22,0 m, Radius 4,2 m) mit einem Kegeldach; die Mantellinie des Kegels ist 7,9 m lang, sein Radius ebenfalls 4,2 m. Bestimme die Gesamthöhe des Turms.
\gzwei{\begin{tikzpicture}[line width=0.7pt,font=\small,scale=0.12]\draw (-4.2,0) arc (180:360:4.2 and 1.47);\draw[dashed,line width=0.5pt] (-4.2,0) arc (180:0:4.2 and 1.47);\draw (-4.2,0) -- (-4.2,22) (4.2,0) -- (4.2,22);\draw (-4.2,22) arc (180:360:4.2 and 1.47);\draw[dashed,line width=0.5pt] (-4.2,22) arc (180:0:4.2 and 1.47);\draw (-4.2,22) -- (0,28.69) -- (4.2,22);\draw[dashed,line width=0.5pt] (0,0) -- (4.2,0);\node[anchor=north] at (2.1,-1.6) {$r = 4{,}2$ m};\node[anchor=west] at (4.4,11) {$22{,}0$ m};\node[anchor=south west] at (2.2,25.4) {$7{,}9$ m};\end{tikzpicture}}{\karo{4}}
\fusshilfe{\gfuss{12}{28,7 m}}
\end{pfaufg}
% 2022-OS-K2c-m1 (Original 2022-OS-K2c, Schwere III, rang 10; herausgelöste Fassung: 2022-OS-K2c-h1)
\begin{pfaufg}{nach~P10\\’22~\pkt{}~2c}{13.}{FOR}
Ein Becher hat die Form eines Zylinders ohne Deckel; er ist 7,5 cm hoch und hat den Radius 3,4 cm. Ein Rührstab steht schräg im Becher: Er berührt unten den Rand des Bodens und lehnt oben am gegenüberliegenden Rand. 2 cm des Stabs ragen über den Becher hinaus. Bestimme die Länge des Stabs.
\gzwei{\begin{tikzpicture}[line width=0.7pt,font=\small,scale=0.6]\draw (0,2.8) ellipse (1.28 and 0.45);\draw (-1.28,0) -- (-1.28,2.8);\draw (1.28,0) -- (1.28,2.8);\draw (-1.28,0) arc (180:360:1.28 and 0.45);\draw[dashed,line width=0.5pt] (-1.28,0) arc (180:0:1.28 and 0.45);\draw[line width=1.3pt] (-1.28,0) -- (1.28,2.8) -- ($(-1.28,0)!1.27!(1.28,2.8)$);\draw[dashed,line width=0.5pt] (0,0) -- (1.28,0);\node[anchor=north] at (0.64,-0.5) {$r = 3{,}4$ cm};\node[anchor=west] at (1.32,1.4) {$7{,}5$ cm};\end{tikzpicture}}{\karo{4}}
\fusshilfe{\gfuss{13}{12,1 cm}}
\end{pfaufg}

\gkern{Die Hypotenuse liegt dem rechten Winkel gegenüber; sie ist die längste Seite, die beiden anderen sind die Katheten.\par\smallskip\noindent a² + b² = c²: Die Quadrate über den Katheten sind zusammen so groß wie das Quadrat über der Hypotenuse c.\par\smallskip\noindent Hypotenuse gesucht: c = $\sqrt{a² + b²}$; Kathete gesucht: a = $\sqrt{c² - b²}$.\par\smallskip\noindent In Figur oder Körper: Such den rechten Winkel (eine Höhe steht senkrecht); die Strecken daran sind die Katheten (Höhe, Radius oder Durchmesser), gegenüber liegt die Hypotenuse.}
\gmehr{6}
\end{document}